\subsection{Hyperbola from the polar formula}

Assume
\[
e>1.
\]
Then the denominator
\[
1+e\cos\theta
\]
vanishes at two angles satisfying
\[
\cos\theta=-\frac1e.
\]
These limiting directions are the asymptotic directions of the branch described
from the chosen focus.

For $r\geq0$, the formula
\[
\boxed{r=\frac{p}{1+e\cos\theta}}
\]
traces the branch lying on the same side of the directrix as the chosen focus.
It does not silently represent the entire two-branch hyperbola.

\begin{center}
\begin{tikzpicture}[scale=0.82,>=stealth]
    \draw[axis,->] (-1.8,0)--(7.2,0) node[right] {polar axis};
    \draw[axis,->] (0,-4.2)--(0,4.4);
    \draw[helper,dashed] (0,0)--(6.4,3.7);
    \draw[helper,dashed] (0,0)--(6.4,-3.7);
    \draw[subset,domain=1.45:6.5,samples=130]
        plot (\x,{1.05*sqrt(\x*\x/2.1-1)});
    \draw[subset,domain=1.45:6.5,samples=130]
        plot (\x,{-1.05*sqrt(\x*\x/2.1-1)});
    \coordinate (F) at (0,0);
    \coordinate (P) at (3.0,1.72);
    \draw[blue!70!black,line width=1.1pt,->] (F)--(P)
        node[midway,above left] {$r$};
    \fill[geometricSubsetColor] (F) circle (2.2pt) node[below left] {chosen focus};
    \fill[geometricSubsetColor] (P) circle (2.2pt) node[above right] {$P$};
    \node at (5.0,-3.7) {one focus-centred branch};
\end{tikzpicture}
\end{center}

To recover the complete hyperbola, return to the two-focus geometry. The second
branch is obtained by choosing the other focus and the corresponding directrix,
or by using the Cartesian equation
\[
\boxed{\frac{x^2}{a^2}-\frac{y^2}{b^2}=1},
\]
which displays both branches at once.

\begin{center}
\begin{tikzpicture}[scale=0.66,>=stealth]
    \draw[axis,->] (-6.2,0)--(6.4,0) node[right] {$x$};
    \draw[axis,->] (0,-4.0)--(0,4.2) node[above] {$y$};
    \draw[helper,dashed] (-5.7,-3.8)--(5.7,3.8);
    \draw[helper,dashed] (-5.7,3.8)--(5.7,-3.8);
    \draw[subset,domain=2.05:5.7,samples=120] plot (\x,{1.333*sqrt(\x*\x/4-1)});
    \draw[subset,domain=2.05:5.7,samples=120] plot (\x,{-1.333*sqrt(\x*\x/4-1)});
    \draw[subset,domain=-5.7:-2.05,samples=120] plot (\x,{1.333*sqrt(\x*\x/4-1)});
    \draw[subset,domain=-5.7:-2.05,samples=120] plot (\x,{-1.333*sqrt(\x*\x/4-1)});
    \fill[geometricSubsetColor] (-3.33,0) circle (2.1pt) node[below] {$F_1$};
    \fill[geometricSubsetColor] (3.33,0) circle (2.1pt) node[below] {$F_2$};
    \node at (0,-4.6) {the complete hyperbola has two branches};
\end{tikzpicture}
\end{center}

\begin{corebox}[Why the polar and Cartesian pictures differ]
The polar equation is adapted to one chosen focus and one directrix. The
Cartesian equation is adapted to the centre and displays both branches
symmetrically. Neither description is incomplete when its geometric viewpoint
is stated correctly.
\end{corebox}
